\IfFileExists{stacks-project.cls}{%
\documentclass{stacks-project}
}{%
\documentclass{amsart}
}

% For dealing with references we use the comment environment
\usepackage{verbatim}
\newenvironment{reference}{\comment}{\endcomment}
%\newenvironment{reference}{}{}
\newenvironment{slogan}{\comment}{\endcomment}
\newenvironment{history}{\comment}{\endcomment}

% For commutative diagrams we use Xy-pic
\usepackage[all]{xy}

% We use 2cell for 2-commutative diagrams.
\xyoption{2cell}
\UseAllTwocells

% We use multicol for the list of chapters between chapters
\usepackage{multicol}

% This is generally recommended for better output
\usepackage{lmodern}
\usepackage[T1]{fontenc}

% For cross-file-references
\usepackage{xr-hyper}

% Package for hypertext links:
\usepackage{hyperref}

% For any local file, say "hello.tex" you want to link to please
% use \externaldocument[hello-]{hello}
\externaldocument[introduction-]{introduction}
\externaldocument[conventions-]{conventions}
\externaldocument[sets-]{sets}
\externaldocument[categories-]{categories}
\externaldocument[topology-]{topology}
\externaldocument[sheaves-]{sheaves}
\externaldocument[sites-]{sites}
\externaldocument[stacks-]{stacks}
\externaldocument[fields-]{fields}
\externaldocument[algebra-]{algebra}
\externaldocument[brauer-]{brauer}
\externaldocument[homology-]{homology}
\externaldocument[derived-]{derived}
\externaldocument[simplicial-]{simplicial}
\externaldocument[more-algebra-]{more-algebra}
\externaldocument[smoothing-]{smoothing}
\externaldocument[modules-]{modules}
\externaldocument[sites-modules-]{sites-modules}
\externaldocument[injectives-]{injectives}
\externaldocument[cohomology-]{cohomology}
\externaldocument[sites-cohomology-]{sites-cohomology}
\externaldocument[dga-]{dga}
\externaldocument[dpa-]{dpa}
\externaldocument[sdga-]{sdga}
\externaldocument[hypercovering-]{hypercovering}
\externaldocument[schemes-]{schemes}
\externaldocument[constructions-]{constructions}
\externaldocument[properties-]{properties}
\externaldocument[morphisms-]{morphisms}
\externaldocument[coherent-]{coherent}
\externaldocument[divisors-]{divisors}
\externaldocument[limits-]{limits}
\externaldocument[varieties-]{varieties}
\externaldocument[topologies-]{topologies}
\externaldocument[descent-]{descent}
\externaldocument[perfect-]{perfect}
\externaldocument[more-morphisms-]{more-morphisms}
\externaldocument[flat-]{flat}
\externaldocument[groupoids-]{groupoids}
\externaldocument[more-groupoids-]{more-groupoids}
\externaldocument[etale-]{etale}
\externaldocument[chow-]{chow}
\externaldocument[intersection-]{intersection}
\externaldocument[pic-]{pic}
\externaldocument[weil-]{weil}
\externaldocument[adequate-]{adequate}
\externaldocument[dualizing-]{dualizing}
\externaldocument[duality-]{duality}
\externaldocument[discriminant-]{discriminant}
\externaldocument[derham-]{derham}
\externaldocument[local-cohomology-]{local-cohomology}
\externaldocument[algebraization-]{algebraization}
\externaldocument[curves-]{curves}
\externaldocument[resolve-]{resolve}
\externaldocument[models-]{models}
\externaldocument[functors-]{functors}
\externaldocument[equiv-]{equiv}
\externaldocument[pione-]{pione}
\externaldocument[etale-cohomology-]{etale-cohomology}
\externaldocument[proetale-]{proetale}
\externaldocument[relative-cycles-]{relative-cycles}
\externaldocument[more-etale-]{more-etale}
\externaldocument[trace-]{trace}
\externaldocument[crystalline-]{crystalline}
\externaldocument[spaces-]{spaces}
\externaldocument[spaces-properties-]{spaces-properties}
\externaldocument[spaces-morphisms-]{spaces-morphisms}
\externaldocument[decent-spaces-]{decent-spaces}
\externaldocument[spaces-cohomology-]{spaces-cohomology}
\externaldocument[spaces-limits-]{spaces-limits}
\externaldocument[spaces-divisors-]{spaces-divisors}
\externaldocument[spaces-over-fields-]{spaces-over-fields}
\externaldocument[spaces-topologies-]{spaces-topologies}
\externaldocument[spaces-descent-]{spaces-descent}
\externaldocument[spaces-perfect-]{spaces-perfect}
\externaldocument[spaces-more-morphisms-]{spaces-more-morphisms}
\externaldocument[spaces-flat-]{spaces-flat}
\externaldocument[spaces-groupoids-]{spaces-groupoids}
\externaldocument[spaces-more-groupoids-]{spaces-more-groupoids}
\externaldocument[bootstrap-]{bootstrap}
\externaldocument[spaces-pushouts-]{spaces-pushouts}
\externaldocument[spaces-chow-]{spaces-chow}
\externaldocument[groupoids-quotients-]{groupoids-quotients}
\externaldocument[spaces-more-cohomology-]{spaces-more-cohomology}
\externaldocument[spaces-simplicial-]{spaces-simplicial}
\externaldocument[spaces-duality-]{spaces-duality}
\externaldocument[formal-spaces-]{formal-spaces}
\externaldocument[restricted-]{restricted}
\externaldocument[spaces-resolve-]{spaces-resolve}
\externaldocument[formal-defos-]{formal-defos}
\externaldocument[defos-]{defos}
\externaldocument[cotangent-]{cotangent}
\externaldocument[examples-defos-]{examples-defos}
\externaldocument[algebraic-]{algebraic}
\externaldocument[examples-stacks-]{examples-stacks}
\externaldocument[stacks-sheaves-]{stacks-sheaves}
\externaldocument[criteria-]{criteria}
\externaldocument[artin-]{artin}
\externaldocument[quot-]{quot}
\externaldocument[stacks-properties-]{stacks-properties}
\externaldocument[stacks-morphisms-]{stacks-morphisms}
\externaldocument[stacks-limits-]{stacks-limits}
\externaldocument[stacks-cohomology-]{stacks-cohomology}
\externaldocument[stacks-perfect-]{stacks-perfect}
\externaldocument[stacks-introduction-]{stacks-introduction}
\externaldocument[stacks-more-morphisms-]{stacks-more-morphisms}
\externaldocument[stacks-geometry-]{stacks-geometry}
\externaldocument[moduli-]{moduli}
\externaldocument[moduli-curves-]{moduli-curves}
\externaldocument[examples-]{examples}
\externaldocument[exercises-]{exercises}
\externaldocument[guide-]{guide}
\externaldocument[desirables-]{desirables}
\externaldocument[coding-]{coding}
\externaldocument[obsolete-]{obsolete}
\externaldocument[fdl-]{fdl}
\externaldocument[index-]{index}

% Theorem environments.
%
\theoremstyle{plain}
\newtheorem{theorem}[subsection]{Theorem}
\newtheorem{proposition}[subsection]{Proposition}
\newtheorem{lemma}[subsection]{Lemma}

\theoremstyle{definition}
\newtheorem{definition}[subsection]{Definition}
\newtheorem{example}[subsection]{Example}
\newtheorem{exercise}[subsection]{Exercise}
\newtheorem{situation}[subsection]{Situation}

\theoremstyle{remark}
\newtheorem{remark}[subsection]{Remark}
\newtheorem{remarks}[subsection]{Remarks}

\numberwithin{equation}{subsection}

% Macros
%
\def\lim{\mathop{\mathrm{lim}}\nolimits}
\def\colim{\mathop{\mathrm{colim}}\nolimits}
\def\Spec{\mathop{\mathrm{Spec}}}
\def\Hom{\mathop{\mathrm{Hom}}\nolimits}
\def\Ext{\mathop{\mathrm{Ext}}\nolimits}
\def\SheafHom{\mathop{\mathcal{H}\!\mathit{om}}\nolimits}
\def\SheafExt{\mathop{\mathcal{E}\!\mathit{xt}}\nolimits}
\def\Sch{\mathit{Sch}}
\def\Mor{\mathop{\mathrm{Mor}}\nolimits}
\def\Ob{\mathop{\mathrm{Ob}}\nolimits}
\def\Sh{\mathop{\mathit{Sh}}\nolimits}
\def\NL{\mathop{N\!L}\nolimits}
\def\CH{\mathop{\mathrm{CH}}\nolimits}
\def\proetale{{pro\text{-}\acute{e}tale}}
\def\etale{{\acute{e}tale}}
\def\QCoh{\mathit{QCoh}}
\def\Ker{\mathop{\mathrm{Ker}}}
\def\Im{\mathop{\mathrm{Im}}}
\def\Coker{\mathop{\mathrm{Coker}}}
\def\Coim{\mathop{\mathrm{Coim}}}

% Boxtimes
%
\DeclareMathSymbol{\boxtimes}{\mathbin}{AMSa}{"02}

%
% Macros for moduli stacks/spaces
%
\def\QCohstack{\mathcal{QC}\!\mathit{oh}}
\def\Cohstack{\mathcal{C}\!\mathit{oh}}
\def\Spacesstack{\mathcal{S}\!\mathit{paces}}
\def\Quotfunctor{\mathrm{Quot}}
\def\Hilbfunctor{\mathrm{Hilb}}
\def\Curvesstack{\mathcal{C}\!\mathit{urves}}
\def\Polarizedstack{\mathcal{P}\!\mathit{olarized}}
\def\Complexesstack{\mathcal{C}\!\mathit{omplexes}}
% \Pic is the operator that assigns to X its picard group, usage \Pic(X)
% \Picardstack_{X/B} denotes the Picard stack of X over B
% \Picardfunctor_{X/B} denotes the Picard functor of X over B
\def\Pic{\mathop{\mathrm{Pic}}\nolimits}
\def\Picardstack{\mathcal{P}\!\mathit{ic}}
\def\Picardfunctor{\mathrm{Pic}}
\def\Deformationcategory{\mathcal{D}\!\mathit{ef}}
\usepackage{cleveref}
\crefname{originalreference}{original reference}{original references}
\Crefname{originalreference}{Original reference}{Original references}
\makeatletter
\newlabel{bootstrap-theorem-final-bootstrap}{{04S6}{source}{}{sourcepointer-1}{}}
\newlabel{bootstrap-theorem-final-bootstrap@cref}{{[originalreference][0][]04S6}{[1][0][]source}{}{}{}}
\makeatother
\crefname{originallemma}{original lemma}{original lemmas}
\Crefname{originallemma}{Original lemma}{Original lemmas}
\makeatletter
\newlabel{lemma-surjection-space-of-sections}{{05XQ}{source}{}{sourcepointer-2}{}}
\newlabel{lemma-surjection-space-of-sections@cref}{{[originallemma][0][]05XQ}{[1][0][]source}{}{}{}}
\makeatother
\crefname{originalreference}{original reference}{original references}
\Crefname{originalreference}{Original reference}{Original references}
\makeatletter
\newlabel{spaces-more-groupoids-proposition-finite-algebraic-space}{{04QH}{source}{}{sourcepointer-3}{}}
\newlabel{spaces-more-groupoids-proposition-finite-algebraic-space@cref}{{[originalreference][0][]04QH}{[1][0][]source}{}{}{}}
\makeatother
\crefname{originalreference}{original reference}{original references}
\Crefname{originalreference}{Original reference}{Original references}
\makeatletter
\newlabel{spaces-more-groupoids-section-finite}{{04PB}{source}{}{sourcepointer-4}{}}
\newlabel{spaces-more-groupoids-section-finite@cref}{{[originalreference][0][]04PB}{[1][0][]source}{}{}{}}
\makeatother
\crefname{originalreference}{original reference}{original references}
\Crefname{originalreference}{Original reference}{Original references}
\makeatletter
\newlabel{spaces-more-morphisms-lemma-where-isomorphism}{{05XD}{source}{}{sourcepointer-5}{}}
\newlabel{spaces-more-morphisms-lemma-where-isomorphism@cref}{{[originalreference][0][]05XD}{[1][0][]source}{}{}{}}
\makeatother
\crefname{originalreference}{original reference}{original references}
\Crefname{originalreference}{Original reference}{Original references}
\makeatletter
\newlabel{spaces-morphisms-lemma-etale-universally-injective-open}{{05W5}{source}{}{sourcepointer-6}{}}
\newlabel{spaces-morphisms-lemma-etale-universally-injective-open@cref}{{[originalreference][0][]05W5}{[1][0][]source}{}{}{}}
\makeatother
\crefname{originalreference}{original reference}{original references}
\Crefname{originalreference}{Original reference}{Original references}
\makeatletter
\newlabel{spaces-morphisms-lemma-section-immersion}{{03KP}{source}{}{sourcepointer-7}{}}
\newlabel{spaces-morphisms-lemma-section-immersion@cref}{{[originalreference][0][]03KP}{[1][0][]source}{}{}{}}
\makeatother
\crefname{originalreference}{original reference}{original references}
\Crefname{originalreference}{Original reference}{Original references}
\makeatletter
\newlabel{spaces-properties-lemma-etale-permanence}{{03FV}{source}{}{sourcepointer-8}{}}
\newlabel{spaces-properties-lemma-etale-permanence@cref}{{[originalreference][0][]03FV}{[1][0][]source}{}{}{}}
\makeatother
\title{The section functor of an etale morphism over a finite locally free base is an etale algebraic space}
\author{The Stacks Project Authors}
\date{stacks-project:criteria.tex:local-snapshot-20261012; GNU FDL 1.2 or later}
\begin{document}
\maketitle
\noindent\textit{Editorial note.} No omission markers in-span; proof complete as written.\par\medskip
\noindent\textit{Allowed standard prerequisites (editorial).}\begin{itemize}
\item Sections are immersions. spaces-morphisms-lemma-section-immersion (cross-chapter, named). Shows sigma is a closed immersion.
\item Etale permanence and open immersion. spaces-properties-lemma-etale-permanence, spaces-morphisms-lemma-etale-universally-injective-open (cross-chapter, named). Upgrade sigma to an open immersion.
\item Finite part is algebraic. spaces-more-groupoids-proposition-finite-algebraic-space (cross-chapter, named). Gives that (W/U)\_fin is an etale algebraic space.
\item Where-isomorphism locus. spaces-more-morphisms-lemma-where-isomorphism (cross-chapter, named). Identifies F as an open of the finite part.
\item Bootstrap. bootstrap-theorem-final-bootstrap (cross-chapter, named). Concludes algebraicity of F in the general case.
\end{itemize}\medskip
\ifdefined\hypertarget\hypertarget{sourcepointer-1}{}\fi
\noindent\textit{Source reference.} Original reference 04S6: focused-seeds/geometry/S9-crit-space-of-sections/ref-tags.txt, source lines 5--5. Navigation pointer only: the counted proof cites this statement by \textbackslash{}ref as a standing input; the deduction the counted proof has to make is carried out inside the counted span, and no part of that deduction is delegated to this pointer.\par\medskip
\ifdefined\hypertarget\hypertarget{sourcepointer-2}{}\fi
\noindent\textit{Source reference.} Original lemma 05XQ: sources/stacks-criteria/pinned-original-source.tex, source lines 1321--1321. Navigation pointer only: the counted proof cites this statement by \textbackslash{}ref as a standing input; the deduction the counted proof has to make is carried out inside the counted span, and no part of that deduction is delegated to this pointer.\par\medskip
\ifdefined\hypertarget\hypertarget{sourcepointer-3}{}\fi
\noindent\textit{Source reference.} Original reference 04QH: focused-seeds/geometry/S9-crit-space-of-sections/ref-tags.txt, source lines 6--6. Navigation pointer only: the counted proof cites this statement by \textbackslash{}ref as a standing input; the deduction the counted proof has to make is carried out inside the counted span, and no part of that deduction is delegated to this pointer.\par\medskip
\ifdefined\hypertarget\hypertarget{sourcepointer-4}{}\fi
\noindent\textit{Source reference.} Original reference 04PB: focused-seeds/geometry/S9-crit-space-of-sections/ref-tags.txt, source lines 7--7. Navigation pointer only: the counted proof cites this statement by \textbackslash{}ref as a standing input; the deduction the counted proof has to make is carried out inside the counted span, and no part of that deduction is delegated to this pointer.\par\medskip
\ifdefined\hypertarget\hypertarget{sourcepointer-5}{}\fi
\noindent\textit{Source reference.} Original reference 05XD: focused-seeds/geometry/S9-crit-space-of-sections/ref-tags.txt, source lines 8--8. Navigation pointer only: the counted proof cites this statement by \textbackslash{}ref as a standing input; the deduction the counted proof has to make is carried out inside the counted span, and no part of that deduction is delegated to this pointer.\par\medskip
\ifdefined\hypertarget\hypertarget{sourcepointer-6}{}\fi
\noindent\textit{Source reference.} Original reference 05W5: focused-seeds/geometry/S9-crit-space-of-sections/ref-tags.txt, source lines 9--9. Navigation pointer only: the counted proof cites this statement by \textbackslash{}ref as a standing input; the deduction the counted proof has to make is carried out inside the counted span, and no part of that deduction is delegated to this pointer.\par\medskip
\ifdefined\hypertarget\hypertarget{sourcepointer-7}{}\fi
\noindent\textit{Source reference.} Original reference 03KP: focused-seeds/geometry/S9-crit-space-of-sections/ref-tags.txt, source lines 10--10. Navigation pointer only: the counted proof cites this statement by \textbackslash{}ref as a standing input; the deduction the counted proof has to make is carried out inside the counted span, and no part of that deduction is delegated to this pointer.\par\medskip
\ifdefined\hypertarget\hypertarget{sourcepointer-8}{}\fi
\noindent\textit{Source reference.} Original reference 03FV: focused-seeds/geometry/S9-crit-space-of-sections/ref-tags.txt, source lines 11--11. Navigation pointer only: the counted proof cites this statement by \textbackslash{}ref as a standing input; the deduction the counted proof has to make is carried out inside the counted span, and no part of that deduction is delegated to this pointer.\par\medskip
\begin{lemma}
\label{lemma-space-of-sections}
Let $Z \to U$ be a finite locally free morphism of schemes.
Let $W$ be an algebraic space and let $W \to Z$ be an \'etale morphism.
Then the functor
$$
F : (\Sch/U)_{fppf}^{opp} \longrightarrow \textit{Sets},
$$
defined by the rule
$$
U' \longmapsto
F(U') =
\{\sigma : Z_{U'} \to W_{U'}\text{ section of }W_{U'} \to Z_{U'}\}
$$
is an algebraic space and the morphism $F \to U$ is \'etale.
\end{lemma}

\begin{proof}
Assume first that $W \to Z$ is also separated.
Let $U'$ be a scheme over $U$ and let $\sigma \in F(U')$. By
Morphisms of Spaces, Lemma \ref{spaces-morphisms-lemma-section-immersion}
the morphism $\sigma$ is a closed immersion.
Moreover, $\sigma$ is \'etale by
Properties of Spaces, Lemma \ref{spaces-properties-lemma-etale-permanence}.
Hence $\sigma$ is also an open immersion, see
Morphisms of Spaces,
Lemma \ref{spaces-morphisms-lemma-etale-universally-injective-open}.
In other words, $Z_\sigma = \sigma(Z_{U'}) \subset W_{U'}$ is
an open subspace such that the morphism $Z_\sigma \to Z_{U'}$
is an isomorphism. In particular, the morphism $Z_\sigma \to U'$
is finite. Hence we obtain a transformation of functors
$$
F \longrightarrow (W/U)_{fin}, \quad
\sigma \longmapsto (U' \to U, Z_\sigma)
$$
where $(W/U)_{fin}$ is the finite part of the morphism $W \to U$
introduced in
More on Groupoids in Spaces, Section
\ref{spaces-more-groupoids-section-finite}.
It is clear that this transformation of functors is injective (since we can
recover $\sigma$ from $Z_\sigma$ as the inverse of the isomorphism
$Z_\sigma \to Z_{U'}$). By
More on Groupoids in Spaces, Proposition
\ref{spaces-more-groupoids-proposition-finite-algebraic-space}
we know that $(W/U)_{fin}$ is an algebraic space \'etale over $U$.
Hence to finish the proof in this case it suffices to show that
$F \to (W/U)_{fin}$ is representable and an open immersion.
To see this suppose that we are given a morphism of schemes $U' \to U$
and an open subspace $Z' \subset W_{U'}$ such that $Z' \to U'$
is finite. Then it suffices to show that there exists an
open subscheme $U'' \subset U'$ such that a morphism
$T \to U'$ factors through $U''$ if and only if $Z' \times_{U'} T$
maps isomorphically to $Z \times_{U'} T$. This follows from
More on Morphisms of Spaces, Lemma
\ref{spaces-more-morphisms-lemma-where-isomorphism}
(here we use that $Z \to B$ is flat and locally of finite presentation
as well as finite).
Hence we have proved the lemma in case $W \to Z$ is separated
as well as \'etale.

\medskip\noindent
In the general case we choose a separated scheme $W'$ and a surjective
\'etale morphism $W' \to W$. Note that the morphisms $W' \to W$ and
$W \to Z$ are separated as their source is separated. Denote $F'$ the
functor associated to $W' \to Z \to U$ as in the lemma. In the first
paragraph of the proof we showed that $F'$ is representable by an
algebraic space \'etale over $U$. By
Lemma \ref{lemma-surjection-space-of-sections}
the map of functors $F' \to F$ is surjective for the \'etale topology
on $\Sch/U$. Moreover, if $U'$ and $\sigma : Z_{U'} \to W_{U'}$
define a point $\xi \in F(U')$, then the fibre product
$$
F'' = F' \times_{F, \xi} U'
$$
is the functor on $\Sch/U'$ associated to the morphisms
$$
W'_{U'} \times_{W_{U'}, \sigma} Z_{U'} \to Z_{U'} \to U'.
$$
Since the first morphism is separated as a base change of a separated
morphism, we see that $F''$ is an algebraic space \'etale over $U'$
by the result of the first paragraph. It follows that $F' \to F$ is a
surjective \'etale transformation of functors, which is representable
by algebraic spaces. Hence $F$ is an algebraic space by
Bootstrap, Theorem \ref{bootstrap-theorem-final-bootstrap}.
Since $F' \to F$ is an \'etale surjective morphism of algebraic spaces
it follows that $F \to U$ is \'etale because $F' \to U$ is \'etale.
\end{proof}

\end{document}
